% 第2章 2.5.3：正交投影与斜投影；偶、奇部分。
% (a) R^2 中 x=(0.6,0.9)：正交投影到横轴得 (0.6,0)；斜投影 P(x_1,x_2)=(x_1+x_2,0) 得 (1.5,0)，
%     其核为直线 x_1+x_2=0，残差 x-Px=(-0.9,0.9) 沿该方向，不与横轴正交；‖P‖=sqrt2。比例 1 单位=20mm。
% (b) u(x)=e^x 在 [-1,1] 上的偶部分 cosh x 与奇部分 sinh x（解析式）。
\begin{tikzpicture}[font=\figurefont\footnotesize,text=FigureInk,
  note/.style={font=\FigureNoteFont\fontsize{7.5}{9}\selectfont,text=FigureMuted},
  vec/.style={line width=1pt,-{Stealth[length=2mm]}}]
\begin{scope}[x=20mm,y=20mm]
  \draw[FigureStroke,line width=.6pt,-{Stealth[length=1.6mm]}] (-0.45,0)--(1.85,0) node[below left=-.2mm] {$M$};
  \draw[FigureStroke,line width=.6pt,-{Stealth[length=1.6mm]}] (0,-0.3)--(0,1.15);
  \draw[FigureGrid,line width=.8pt] (-0.4,0.4)--(0.3,-0.3);
  \node[note,anchor=north] at (-0.35,0.25) {$\ker P$};
  \draw[vec,FigureInk] (0,0)--(0.6,0.9) node[above] {$\bm x$};
  \draw[densely dashed,FigureBlue,line width=.8pt] (0.6,0.9)--(0.6,0);
  \draw[densely dashed,FigureYellow,line width=.8pt] (0.6,0.9)--(1.5,0);
  \fill[FigureBlue] (0.6,0) circle (1.4pt) node[below=.6mm,text=FigureBlue] {$P_M\bm x$};
  \fill[FigureYellow] (1.5,0) circle (1.4pt) node[below=.6mm,text=FigureYellow!70!FigureInk] {$P\bm x$};
  \draw[FigureBlue,line width=.5pt] (0.6,0.08)--(0.68,0.08)--(0.68,0);
  \node[anchor=west] at (-0.5,1.35) {(a) 正交投影与斜投影};
\end{scope}
\begin{axis}[at={(48mm,-6mm)},anchor=south west,width=56mm,height=38mm,axis lines=left,
  x axis line style={draw=FigureStroke,line width=.6pt},y axis line style={draw=FigureStroke,line width=.6pt},
  tick style={draw=FigureStroke},tick label style={font=\figurefont\scriptsize,text=FigureInk},
  xlabel={$x$},xlabel style={font=\figurefont\scriptsize,at={(axis description cs:1,0)},anchor=north west,yshift=1mm},
  xmin=-1,xmax=1,ymin=-1.3,ymax=2.9,xtick={-1,0,1},ytick={-1,0,1,2},
  ymajorgrids,grid style={draw=FigureGrid,line width=.4pt},samples=80,domain=-1:1,
  title={(b) $e^x=\cosh x+\sinh x$},title style={font=\figurefont\footnotesize},
  legend style={draw=none,fill=none,font=\figurefont\scriptsize,cells={anchor=west},at={(0.5,-0.22)},anchor=north},legend columns=3]
\addplot[draw=FigureInk,line width=1.1pt] {exp(x)};
\addplot[draw=FigureBlue,line width=1.1pt,dashed] {cosh(x)};
\addplot[draw=FigureGreen,line width=1.1pt,dash dot] {sinh(x)};
\legend{$e^x$,{$\cosh x$（偶）},{$\sinh x$（奇）}}
\end{axis}
\end{tikzpicture}
